%%%
%
% $Autor: Gruppe 05 $
% $Datum: 2025-05-02 $
% $Dateiname: Pumpe
% $Version: 1 $
%
% !TeX spellcheck = de_DE/GB
% !TeX program = pdflatex
% !BIB program = biber/bibtex
% !TeX encoding = utf8
%
%%%

\chapter{Comet Tauchpumpe ELEGANT 12V}
\section{Einführung}
Pumpen sind Einrichtungen zur Förderung von flüssigen Medien von einem niedrigeren auf ein höheres statisches Druckniveau. Dieser Vorgang wird über verschiedene Wege erreicht. 
Ein höheres statisches Druckniveau kann durch folgende Verfahren erreicht werden: 
\begin{itemize}
	\item das Einwirken einer Druckraft, 
	\item die Übertragung mechanischer Arbeit, 
	\item den Impulsaustausch, 
	\item das Zumischen von Druckluft zum Förderwasser, 
	\item die Stoßwirkung einer in ihrer Bewegung gehemmten Wassersäule oder 
	\item durch Ausnutzung der Zähigkeit des Mediums unter Verwendung eines Förder-gewindes 
\end{itemize}
All die verschiedenen Arbeitsweisen beeinflussen die Pumpenbauart. Die häufigsten vorkommenden Pumpen sind die Verdränger- und Kreiselpumpe. 
Im Folgenden wird näher auf die Kreiselpumpe eingegangen. 
Kreispumpen gibt es in ein- oder mehrstufigen Ausführungen. Ihre Arbeitsweise ist gleich, sie unterscheiden sich lediglich in ihrem Aufbau.\cite{Schulz:2013}
Pumpen sind oft Komponenten von komplexen Prozessen, deren Ausfall erhebliche Auswirkungen verursachen kann. Die Anforderung an maximale Betriebssicherheit ist daher von besonderer Bedeutung.\cite{Gulich:2020}

\subsection{Arbeitsweise Kreiselpumpen}
Die Kreiselpumpe verfügt über ein rotierendes Laufrad mit Schaufeln. Dieses Rad über-trägt mechanische Arbeit auf die in den Laufradkanälen befindliche Flüssigkeit, die da-bei durch Fliehkräfte aus dem Laufrad verdrängt wird. Durch den Vorgang kommt es zur Druckerhöhung und Geschwindigkeitszunahme des Mediums. Die gleichzeitige Zunahme der Absolutgeschwindigkeit des Fördermittels stellt eine unerwünschte Strömungsbegleiterscheinung dar, da in der Pumpe vorrangig eine Druckerhöhung erzielt werden soll. Daher muss diese anschließend in Druckenergie umgewandelt werden. Die Umwandlung wird durch das um das Laufrad befindliche ringförmige Leitrad realisiert. Zusammen bilden Lauf- und Leitrad eine Stufe (daher einstufige Kreiselpumpe).
Die kontinuierliche Aufrechterhaltung der Strömung erfolgt dadurch, dass durch die aus dem drehenden Laufrad verdrängte Flüssigkeit eine Saugwirkung erzeugt und das gleiche Volumen an Flüssigkeit über Saugstutzen wieder in die Pumpe eintritt. 
Die Wahl einer Kreiselpumpe erfolgt nach Förderhöhe und Druckhöhe. So kann nach der Förderhöhe in Nieder, - Mittel – und Hochdruckpumpen unterteilt werden. Im Hin-blick auf die Druckhöhe, wählt man bei großen Druckhöhen mehrstufige Kreiselpumpen. Sie sind im Aufbau wie die einstufigen Kreiselpumpen mit dem einzigen Unter-schied, dass mehrere Stufen (gleichartige Lauf- und Leiträder) hintereinandergeschaltet sind. Dabei wird die Anzahl der Stufen durch die Förderhöhe, Drehzahl der Pumpe und die Größe des Volumenstromes bestimmt. 
Bei geringen Förderhöhen und großen Volumenströmen entspricht die Laufradform, die dem Francisrade bei Wasserturbinen. Die Schaufeln des Laufrades sind bei dieser Bauart räumlich gekrümmt bzw. verwunden. Eine Erhöhung des Volumenstroms bei gering bleibender Förderhöhe erfordert eine mehrströmige Ausführung der Pumpe. Bei dieser Bauart sind mehrere Laufräder parallelgeschaltet. Verwirklicht wird das Ganze durch zwei Einzelräder, welche zu einem Radpaar mit doppelseitigem Wassereintritt vereinigt werden. Beispiele für mehrströmige Pumpen sind Kühlwasserpumpen für Kondensationsanlagen. Eine weitere Anwendung bei gleichbleibenden Bedingungen (Vergrößerung Volumenstrom, klein bleibende Förderhöhe) findet die Pumpe halbaxialer Bauart. Diese Bauart zeichnet sich dadurch aus, dass die axial eintretende Flüssigkeit nur noch eine geringe Umlenkung nach der radialen Richtung erfährt. Sie ersetzt die mehrstufige und bildet den Übergang von der Radial- zur Axialpumpe. Reine Axialpumpen eignen sich besonders zur Überwindung großer Volumenströme bei geringen Strecken.\cite{Schulz:2013}


